\section{Analyticity of the boundary}\label{sec:analytic}

The phase graphs give exact equations for the boundary, one at each
sufficiently high Fourier frequency.  The decisive analytic fact is
that the derivative of these equations is close to the identity on
the high modes.  We state the two estimates imported from
\cite[Lemma~2.3 and Proposition~2.4]{NB}, then prove the regularity
conclusion from our phase hypothesis. The grazing coordinate follows
the normalization introduced by \citet{Lazutkin1973}; asymptotic grazing
actions and rational Fourier tests also appear, in different settings,
in \citet{MarviziMelrose1982}, \citet{RamirezRos2006}, and \citet{ADK2016}.

\begin{theorem}\label{thm:analytic-boundary}
Under the hypotheses of Theorem~\ref{thm:main}, the support function
\(h\) and the boundary of the table are real analytic.
\end{theorem}

\subsection{Normalization and the two analytic inputs}

Choose \(k>0\) so that \(\dd x/\dd\theta=k r(\theta)^{1/3}\)
defines a coordinate of period \(2\pi\).  With \(\theta=\phi(x)\), put
\begin{equation}\label{ana:normalization}
 u=\log\phi',\qquad c_*=k^{-3},\qquad
 r(\phi(x))=c_*e^{-3u(x)}.
\end{equation}
The normalized action is
\[
 D_u(X,Y)=c_*^{-1}\left\{
 \cS(\phi(X),\phi(Y))-
 \int_{\phi(X)}^{\phi(Y)}h(\eta)\,\dd\eta\right\}.
\]
The endpoint integral telescopes along a chain, so the Euler equations
are unchanged by its subtraction and the increasing coordinate change.
Writing \(\phi(\mu)=(\phi(X)+\phi(Y))/2\), integration by parts gives
\begin{equation}\label{eq:normalized-action}
\begin{split}
D_u(X,Y)={}&-\int_X^\mu
 \bigl(1-\cos(\phi(z)-\phi(X))\bigr)e^{-2u(z)}\,\dd z\\
 &-\int_\mu^Y
 \bigl(1-\cos(\phi(Y)-\phi(z))\bigr)e^{-2u(z)}\,\dd z.
\end{split}
\end{equation}
Indeed, before the coordinate change, the sum of the two integrals
with integrands \(r(\eta)(1-\cos(\eta-\theta))\) and
\(r(\eta)(1-\cos(\theta'-\eta))\) is
\(\int_\theta^{\theta'}h-\cS(\theta,\theta')\); the terms containing
\(h'\) at the midpoint cancel.  Now use
\(r(\phi(z))\phi'(z)/c_*=e^{-2u(z)}\).

For an arbitrary smooth real periodic function \(f\), define \(D_f\)
by \eqref{eq:normalized-action}, replacing \(u\) by \(f\) and taking
any primitive \(\phi_f\) of \(e^f\).  No condition is imposed on
\(\int_0^{2\pi}e^f\).  Translating both endpoints by \(2\pi\)
adds the same constant to every primitive value, so \(D_f\) is
diagonally periodic.  The midpoint is well defined for small steps
because \(\phi_f'>0\) on the real axis.  The same formula defines
the holomorphic action locally near the diagonal for holomorphic
real-symmetric inputs.  This arbitrary-input formulation is essential:
the trial functions below need not themselves represent billiard tables.

The normalization fixes the leading action independently of the input:
\begin{equation}\label{ana:universal-cubic}
 D_f(x,x+b)=-\frac{b^3}{24}+O(b^5).
\end{equation}
Indeed, the midpoint is \(x+b/2+O(b^2)\), and
\(\phi_f'^2e^{-2f}=1\), so each half-integral has leading value
\(b^3/48\). Reversing the endpoints gives
\(D_f(Y,X)=-D_f(X,Y)\). If \(a_3(x),a_4(x)\) are its cubic and
quartic coefficients, this identity gives \(2a_4=a_3'\); hence the
quartic coefficient vanishes. This common cubic term explains the
uniform mesh used as the reference in the next proposition. The
Fourier-row estimate below separately controls variation with the input.

For \(0<\sigma\le1\), let
\(\mathcal A_\sigma=\{z\in\C:|\operatorname{Im}z|<\sigma\}\) and set
\[
 \|f\|_{j,\sigma}=
 \max_{0\le i\le j}\sup_{\mathcal A_\sigma}|f^{(i)}|.
\]
Inputs on positive strips are periodic, holomorphic, and real symmetric;
at \(\sigma=0\) they are smooth real periodic functions and the norm
is the corresponding real \(C^j\) norm.

\begin{proposition}[Stationary meshes, {\cite[Lemma~2.3]{NB}}]
\label{prop:mesh-input}
Fix an integer \(m\ge2\) and a finite bound \(M\), and put
\(p=m+2\), \(r_0=m+15\).  There are \(R,K,N\), depending only on
\(m,M\), with the following properties.  For
\(0\le\sigma\le1\), \(\|f\|_{r_0,\sigma}\le M\), and \(q\ge N\),
the action \(D_f\) has an interior-stationary mesh with endpoints
\(X_0=x\), \(X_q=x+2\pi\), of the form
\[
 X_j(x)=x+2\pi j/q+w_j(x),\qquad w_0=w_q=0.
\]
On the strip of width \(\sigma'=(1-K/q^2)\sigma\),
\begin{equation}\label{ana:mesh-bounds}
 q^2\max_j\|w_j\|_{p,\sigma'}+
 q^3\max_j\|w_{j+1}-w_j\|_{p,\sigma'}\le R.
\end{equation}
The mesh is real symmetric and differentiable in smooth real input
variations; its derivative in the norm on the left of
\eqref{ana:mesh-bounds} is bounded by
\(C\|\dot f\|_{r_0,\sigma}\).  The corresponding difference bound
holds for two inputs.  On the real axis the mesh is pointwise unique
in the ball with these scaled position and step bounds.  Any fixed
enlargement of that real uniqueness ball is allowed after increasing
\(N\).
\end{proposition}

Define the critical action and its selected Fourier coefficients by
\begin{equation}\label{ana:constraints}
\begin{split}
 W_q(x;f)&=\sum_{j=0}^{q-1}D_f(X_j(x),X_{j+1}(x)),\\
 C_n(f)&=-\frac{q^2}{6\pi}[W_q(\cdot;f)]_n,
 \qquad n=\pm q,\\
 [v]_n&=\frac1{2\pi}\int_0^{2\pi}v(x)e^{-\ii nx}\,\dd x.
\end{split}
\end{equation}
Thus \(W_q\) records how the stationary-chain action depends on its
initial point, and \(C_q,C_{-q}\) select its \(q\)-th Fourier harmonics.
Periodicity of the displacements and critical action follows from real
mesh uniqueness.  The same
uniqueness identifies these definitions on common real inputs when
different strip widths or different fixed smoothness bounds are used,
for all sufficiently large rows.  In comparing two constructions, one
enlarges the real uniqueness ball to contain both meshes.  Thus this
compatibility is independent of a chosen holomorphic extension.

\begin{proposition}[Fourier rows, {\cite[Proposition~2.4]{NB}}]
\label{prop:row-input}
For fixed \(m\ge2\), \(M<\infty\), and \(r_0=m+15\), there are
\(K,N\), independent of \(0\le\sigma\le1\), such that for
\(\|f\|_{r_0,\sigma}\le M\) and \(|n|\ge N\),
\begin{equation}\label{ana:row-derivative}
 (\dd C_n)_f(v)=[b_n(f)v]_n,
\end{equation}
where
\begin{align}
 \|b_n(f)-1\|_{m,(1-K/n^2)\sigma}&\le K/|n|,
 \label{ana:row-bound}\\
 \|b_n(f)-b_n(\widehat f)\|_{m,(1-K/n^2)\sigma}
 &\le (K/|n|)\|f-\widehat f\|_{r_0,\sigma}.
 \label{ana:row-difference}
\end{align}
The second estimate holds for two inputs with the same bound.
The derivative identity holds for smooth real variations and for
real-symmetric holomorphic variations with the indicated finite norms.
\end{proposition}

Both inputs concern the arbitrary-input action
\eqref{eq:normalized-action}; neither presupposes physical caustics.
We use their estimates directly, without applying the analytic-boundary
result of \cite{NB} under its stronger standing geometric hypothesis.

\subsection{The equations satisfied by the actual boundary}

\begin{lemma}\label{ana:vanishing}
For the smooth function \(u\) in \eqref{ana:normalization},
\(C_n(u)=0\) for every sufficiently large \(|n|\).
\end{lemma}

\begin{proof}
Apply Proposition~\ref{prop:mesh-input} on the real axis.
Its steps are \(2\pi/q+O(q^{-3})\), hence positive and uniformly
small.  Applying \(\phi\) gives a stationary chain for \(\cS\),
with positive uniformly small angular increments and total advance
\(2\pi\).  Proposition~\ref{prop:graph-collar} and
Lemma~\ref{lem:stationary-closure} show that its final momentum equals
its initial momentum.  Thus the chain is stationary also at the
identified endpoints.  More explicitly, differentiation of its
critical action cancels every interior term, and the two endpoint
terms are
\[
 W_q'(x;u)=\frac{\phi'(x)}{c_*}(p_q-p_0)=0.
\]
Here \(\phi'(x+2\pi)=\phi'(x)\), and periodicity of \(h\) cancels
the two endpoint contributions from the subtracted integral.  The
action is constant in \(x\), so its nonzero Fourier coefficients vanish.
\end{proof}

\subsection{A bounded inverse on one Fourier strip}

For a Fourier sequence and \(a\ge0\), define
\[
 |v|_{\sigma,a}=\sum_{k\in\mathbb Z}
 (1+|k|)^a|[v]_k|e^{\sigma|k|},\qquad
 |v|_\sigma=|v|_{\sigma,0}.
\]
The norm \(|\cdot|_\sigma\) defines a Wiener space.  Let \(\Pi_H\)
be the projection to \(|k|\ge N\), and write \(C(f)\) for the
sequence of these output rows.

\begin{lemma}\label{ana:wiener}
Fix the bounds in Proposition~\ref{prop:row-input} with \(m=2\).
The row formula extends to a bounded Fourier operator
\[
 A(f)=\Pi_H+R(f),\qquad
 \|R(f)\|\le C/N,
\]
on each Wiener space of width \(\sigma\).  For two bounded inputs,
\[
 \|R(f)-R(\widehat f)\|
 \le C\|f-\widehat f\|_{r_0,\sigma}.
\]
For sufficiently large \(N\), the restriction \(A_H(f)\) to the
high input modes is invertible with \(\|A_H(f)^{-1}\|\le2\).
The operators and this inverse preserve real Fourier symmetry.
With higher fixed input bounds, the same matrix is bounded in
\(|\cdot|_{\sigma,a}\) whenever \(m>a+1\).
\end{lemma}

\begin{proof}
Contour displacement and \(m\) integrations by parts in
\eqref{ana:row-bound} give
\begin{equation}\label{ana:coefficient-decay}
 |[b_n(f)-1]_j|\le
 \frac{C}{|n|}(1+|j|)^{-m}
 e^{-(1-K/n^2)\sigma|j|}.
\end{equation}
For positive \(\sigma\), displace first to an interior closed strip
and then take its width to the indicated limit; for \(\sigma=0\),
only integration by parts is needed.  Coefficient differences satisfy
the same bound multiplied by
\(\|f-\widehat f\|_{r_0,\sigma}\).

Increase \(N\) so that \(N^2\ge2K\), and set
\(\alpha=K/n^2\le1/2\).  The triangle inequality yields the useful
identity and bound
\begin{align}
 |n|-|k|-(1-\alpha)|n-k|
 &=(1-\alpha)(|n|-|k|-|n-k|)
       +\alpha(|n|-|k|)\notag\\
 &\le \alpha|n|=K/|n|.
 \label{ana:exponent-balance}
\end{align}
Consequently the weighted matrix entry from \(k\) to \(n\) of
\(R(f)\) is at most
\[
 \frac{C}{|n|}(1+|n-k|)^{-m}e^{K\sigma/|n|}.
\]
At \(m=2\) its column sums are \(O(N^{-1})\), proving the stated
operator bounds on weighted \(\ell^1\).  A Neumann series gives the
inverse.  The equality \(C_{-q}(f)=\overline{C_q(f)}\) for real
inputs gives the real symmetry.  Finally,
\[
 \frac{(1+|n|)^a}{(1+|k|)^a}\le(1+|n-k|)^a
\]
reduces the polynomially weighted estimate to the convergent sum
of \((1+|n-k|)^{a-m}\).
\end{proof}

The Fourier operator here is a bounded extension of the actual
derivative on sufficiently regular variations.  A Wiener sequence
need not have \(r_0\) bounded derivatives on its original strip.
On every smaller strip it does, and the matrix agrees there with
the actual derivative by Fourier-polynomial approximation.
Compatibility of the real meshes also identifies the row derivatives
across widths.  We will only differentiate on segments with the finite
input bounds required by Proposition~\ref{prop:row-input}.

\subsection{Analytic approximation and real uniqueness}

\begin{proof}[Proof of Theorem~\ref{thm:analytic-boundary}]
Fix \(m=2\), \(r_0=17\), and
\[
 M>4e(1+|u|_{0,r_0}).
\]
Choose \(N\) large enough for Lemmas~\ref{ana:vanishing} and
\ref{ana:wiener}, with \(\|R\|\le1/2\) throughout this input ball.
These choices precede all choices of approximation parameters.

Let \(f_0=\sum_{|k|\le P}[u]_ke^{\ii kx}\), where \(P>N\), and
put \(\sigma_0=P^{-1}\).  For every fixed \(a\ge0\),
\begin{equation}\label{ana:truncation}
 |f_0|_{\sigma_0,a}\le e|u|_{0,a}.
\end{equation}
We claim that, for every fixed positive integer \(a\),
\begin{equation}\label{ana:initial-error}
 |C(f_0)|_{\sigma_0}=O(P^{-a}).
\end{equation}
For \(N\le|n|\le P\), the exponential weights are at most \(e\).
Integrate the real derivative along the segment from \(u\) to \(f_0\),
use \(C(u)=0\), and apply Lemma~\ref{ana:wiener} at width zero.
This part is bounded by \(C|f_0-u|_0\), which decays faster than
any fixed power of \(P\) because \(u\) is smooth.

For \(|n|>P\), instead integrate along \(\lambda f_0\),
\(0\le\lambda\le1\).  The zero input has a uniform stationary mesh
and a constant critical action, so \(C(0)=0\).  For the chosen \(a\),
use Proposition~\ref{prop:row-input} with one fixed integer
\(m>a+1\).  Formula~\eqref{ana:truncation} bounds all its required
input derivatives uniformly in \(P\).  Once \(P\) exceeds this
higher-order row threshold, real mesh uniqueness identifies these
rows with the original ones.  The final assertion of
Lemma~\ref{ana:wiener} gives
\[
 \sum_{|n|>P}(1+|n|)^a|C_n(f_0)|e^{\sigma_0|n|}
 \le C_a|f_0|_{\sigma_0,a}\le C'_a.
\]
Division by \((1+P)^a\) proves the tail bound.  This establishes
\eqref{ana:initial-error}; its constants and threshold may depend
on \(a\), and no uniformity over all smoothness orders is asserted.

Keep the modes \(|k|<N\) fixed and set
\[
 \sigma_j=\frac{1+2^{-j}}{2P},\qquad
 \sigma_\infty=\frac1{2P},\qquad
 \eps_j=|C(f_j)|_{\sigma_j}.
\]
At the \(j\)-th step define
\begin{equation}\label{ana:newton-step}
 \eta_j=-A_H(f_j)^{-1}C(f_j),\qquad f_{j+1}=f_j+\eta_j.
\end{equation}
The correction is real symmetric, has only high modes, and obeys
\(|\eta_j|_{\sigma_j}\le2\eps_j\).  Since
\(\sigma_j-\sigma_{j+1}=2^{-j-2}/P\), the elementary Fourier
smoothing estimate gives
\begin{equation}\label{ana:correction-bound}
 \|\eta_j\|_{r_0,\sigma_{j+1}}
 \le C P^{r_0}2^{r_0j}\eps_j.
\end{equation}
If the correcting segment is in the fixed input ball on this smaller
strip, integration of the derivative difference gives
\begin{equation}\label{ana:newton-bound}
 \eps_{j+1}\le
 C\|\eta_j\|_{r_0,\sigma_{j+1}}|\eta_j|_{\sigma_{j+1}}
 \le C P^{r_0}2^{r_0j}\eps_j^2.
\end{equation}
The linear term cancels by \eqref{ana:newton-step}: compatibility of
the real row derivatives makes the matrix computed at width
\(\sigma_j\) the same matrix on the smaller strip.  Equivalently,
one can integrate each smooth real row first and then take the
smaller-strip norm, avoiding any derivative assertion for rough
variations at the original strip boundary.

Choose \(P\) so large that \(\eps_0\le P^{-2r_0}\), possible by
\eqref{ana:initial-error}, and enlarge it further as needed below.
The simultaneous induction is
\begin{equation}\label{ana:newton-majorant}
 \eps_j\le P^{-2r_0}2^{-(r_0+1)j},
 \qquad
 \|\eta_j\|_{r_0,\sigma_{j+1}}
 \le C P^{-r_0}2^{-j}.
\end{equation}
Indeed, before using \eqref{ana:newton-bound}, the known bound for
\(\eps_j\) and \eqref{ana:correction-bound} put the entire current
segment in the ball: the initial norm is less than \(M/4\), and
the cumulative corrections are at most
\(CP^{-r_0}\sum_{i\ge0}2^{-i}<M/4\).  Thus the quadratic estimate
is applicable.  Its ratio to the proposed bound for \(\eps_{j+1}\)
is at most \(CP^{-r_0}2^{r_0+1-j}\), which is less than one for
one sufficiently large fixed \(P\).  This proves both the error
bound and the input bound at every step.

The summable corrections converge on the positive-width strip
\(\mathcal A_{\sigma_\infty}\), with all derivatives through
\(r_0\), to a real-symmetric holomorphic function \(f_*\).
Its low Fourier modes agree with those of \(u\), and continuity
of each real mesh gives \(C_n(f_*)=0\) for every \(|n|\ge N\).
The real segment between \(u\) and \(f_*\) lies in the same input
ball.  Put \(v=f_*-u=\Pi_Hv\).  Integrating the real derivative,
\[
 0=C(f_*)-C(u)
   =v+\int_0^1R(u+\lambda v)v\,\dd\lambda.
\]
The bound \(\|R\|\le1/2\) implies \(|v|_0\le|v|_0/2\), hence
\(v=0\).  The original smooth \(u\) is therefore analytic.

Finally \(\phi'=e^u>0\), so \(\phi\) and its inverse are analytic.
Equation~\eqref{ana:normalization} makes \(r\) analytic in the normal
angle, and the equation \(h''+h=r\) makes \(h\) analytic.
Thus \(h n_\theta+h'e_\theta\) is an analytic regular boundary
parametrization, proving Theorem~\ref{thm:analytic-boundary}.
\end{proof}

\begin{corollary}\label{ana:glancing-interface}
For the actual table, \(D=D_u\) is real symmetric and holomorphic
in a uniform neighborhood of the real diagonal.  It is diagonally
\(2\pi\)-periodic and antisymmetric, and satisfies
\begin{equation}\label{ana:glancing-expansion}
 D(x,x+b)=-\frac{b^3}{24}+O(b^5),\qquad
 D_1(x,x+b)=\frac{b^2}{8}+O(b^4),
\end{equation}
locally uniformly with each fixed number of derivatives.
\end{corollary}

\begin{proof}
The midpoint implicit function theorem, positivity of \(\phi'\)
on the real axis, and compactness modulo the period give a uniform
holomorphic domain for \eqref{eq:normalized-action}.  Reversing its
two endpoints gives antisymmetry. The calculation of
\eqref{ana:universal-cubic} applies on this holomorphic domain.
Differentiating at fixed second endpoint is the operator
\(\partial_x-\partial_b\), which gives the second expansion.
\end{proof}
